%% ============================================================================
%%  <Course name> --- problem book, English edition.
%%  Build: ./build.sh problems   (latexmk -xelatex, then tools/probcheck.py)
%%
%%  Problems and solutions alternate: each problem starts a new page and its
%%  solution starts the next one (macros in problems/common.tex).  Keep
%%  problems-ko.tex in step, chapter for chapter and problem for problem
%%  (tools/probcheck.py checks that they match).  The web viewer finds the
%%  solution pages by itself, through the named destinations that \problem
%%  and \solution leave in the PDF.
%% ============================================================================
\documentclass[11pt]{article}
\usepackage{coursenotes}

\coursename{Sample Course}
\term{Fall 2027}
\lectitle{Problems and Solutions}
\lecdate{Updated 6 September 2027}

\newcommand{\ProbWord}{Problem}
\newcommand{\SolWord}{Solution}
\newcommand{\SourceWord}{Source}
\newcommand{\NotesWord}{Notes}

%% ============================================================================
%%  problems/common.tex --- shared by problems-en.tex and problems-ko.tex
%%
%%  The problem book alternates problem and solution: every problem starts a
%%  new page, and its solution starts the next new page.  The web viewer
%%  (pdf-notes-viewer, DOCS entry with solutions:true) covers the solution
%%  pages until they are tapped.  It finds them through named destinations
%%  that \problem and \solution put into the PDF:
%%      pb.prob.<n.m>   where problem n.m starts
%%      pb.sol.<n.m>    where its solution starts
%%      pb.end          where back matter after the last solution starts
%%  A solution is covered from its pb.sol page up to the page before the next
%%  mark, or to the last page.  Anything placed after the last solution (an
%%  appendix, say) must therefore start on a new page with \pbmark{end}.
%%  "pb" must equal MARK in the viewer's configuration block.
%%
%%  The labels chapstart:<n>, probstart:<n.m>, solstart:<n.m> are for
%%  tools/probcheck.py, which checks that the two editions match.
%%
%%  Numbering: chapters use the section counter, problems the subsection
%%  counter, so bookmarks read "2 Steady states" and "2.1 A membrane (...)":
%%  the viewer keys on that leading number to keep your place across the
%%  EN/KO switch.  Both editions must carry the same chapters and problems in
%%  the same order.  Equations are numbered afresh in every problem.
%%
%%  Each main file defines the words used below:
%%    \ProbWord \SolWord \SourceWord \NotesWord
%%  Requires the course style (coursenotes.sty: hyperref, enumitem, fancyhdr).
%% ============================================================================

%% ---- notation shared by the problems (add the course's own here) ----------
\newcommand{\ans}[1]{\boxed{\;#1\;}}
\newcommand{\e}{\mathrm{e}}

%% ---- problem parts ---------------------------------------------------------
\newenvironment{parts}%
  {\begin{enumerate}[label=(\alph*),leftmargin=2.2em,itemsep=4pt,topsep=4pt]}%
  {\end{enumerate}}
%% a part of a solution: \solpart{a}{Short title}
\newcommand{\solpart}[2]{%
  \par\addvspace{9pt}\noindent\textbf{(#1)\enspace#2}\par\nobreak\vspace{2pt}}

%% ---- equations restart in every problem ------------------------------------
\makeatletter
\@addtoreset{equation}{subsection}
\renewcommand{\theHequation}{\thesubsection.\arabic{equation}}
\makeatother

%% ---- chapters, problems, solutions ----------------------------------------
\makeatletter
\newif\ifpb@afterchap
\def\pb@probtitle{}

%% \pbmark{name}: a named destination pb.<name> on the current page, for the
%% web viewer (see the top of this file).  Same mechanism as \phantomsection.
\newcommand{\pbmark}[1]{\Hy@raisedlink{\hyper@anchorstart{pb.#1}\hyper@anchorend}}

%% \probchapter{Title}{notes range}
\newcommand{\probchapter}[2]{%
  \clearpage
  \refstepcounter{section}%
  \phantomsection
  \addcontentsline{toc}{section}{\protect\numberline{\thesection}#1}%
  \edef\pb@id{\thesection}%
  \label{chapstart:\pb@id}%
  \markright{#1}%
  \vspace*{-0.2\baselineskip}%
  \noindent\rule{\textwidth}{1.1pt}\par\nobreak\vspace{7pt}%
  \noindent{\Large\bfseries\thesection\quad #1}\par\nobreak\vspace{3pt}%
  \noindent{\small\itshape #2}\par\nobreak\vspace{5pt}%
  \noindent\rule{\textwidth}{0.4pt}\par\nobreak
  \vspace{16pt}%
  \global\pb@afterchaptrue}

%% \problem{Title}{source}{related notes}
\newcommand{\problem}[3]{%
  \ifpb@afterchap\global\pb@afterchapfalse\else\clearpage\fi
  \refstepcounter{subsection}%
  \gdef\pb@probtitle{#1}%
  \phantomsection
  \addcontentsline{toc}{subsection}%
    {\protect\numberline{\thesubsection}#1\texorpdfstring{\ \ {\normalfont\small(#2)}}{ (#2)}}%
  \edef\pb@id{\thesubsection}%
  \label{probstart:\pb@id}%
  \pbmark{prob.\pb@id}%
  \markright{\ProbWord\ \thesubsection}%
  \noindent{\large\bfseries\ProbWord\ \thesubsection\quad #1}\par\nobreak
  \vspace{3pt}%
  \noindent{\small\SourceWord: #2\hfill\NotesWord: #3}\par\nobreak
  \vspace{2pt}\noindent\rule{\textwidth}{0.4pt}\par\nobreak
  \vspace{6pt}}

%% \solution   (starts the solution of the current problem on a new page)
\newcommand{\solution}{%
  \clearpage
  \phantomsection
  \edef\pb@id{\thesubsection}%
  \label{solstart:\pb@id}%
  \pbmark{sol.\pb@id}%
  \markright{\SolWord\ \thesubsection}%
  \noindent{\large\bfseries\SolWord\ \thesubsection\quad\pb@probtitle}\par\nobreak
  \vspace{2pt}\noindent\rule{\textwidth}{0.4pt}\par\nobreak
  \vspace{6pt}}

%% the running head: title on the left, the current problem/solution on the right
\newcommand{\probheads}{%
  \fancyhead[R]{\small\rightmark}}
\makeatother

%% ---- table of contents: problems are listed under their chapter -------------
\setcounter{tocdepth}{2}


\begin{document}
\makelecheader
\probheads

\tableofcontents

\clearpage
\markboth{}{}
\section*{How to use this collection}
Each problem starts on a new page and its solution starts on the next one, so
you can work a problem before you turn the page. In the web viewer the
solution pages stay covered until you tap them.

\emph{Sources.} The source line under each title says where a problem comes
from: a past exam or problem set (with the year), the practice exams written
for these notes, or ``New'' for problems written for this collection. Where a
past problem has extra parts, it says ``+ new parts''.

\emph{Notation} follows the lecture notes; the ``Notes'' field names the
sections a problem rests on.

%% ============================================================================
\probchapter{Random walks and diffusion}{Notes \S1--\S2}

%% ----------------------------------------------------------------------------
\problem{A biased walk}{New}{\S1}

A walker on a line takes independent steps $s_k$, each $+a$ with probability
$p$ and $-a$ with probability $q=1-p$. After $n$ steps it is at
$x_n=\sum_{k=1}^n s_k$.
\begin{parts}
  \item Find the mean and the variance of a single step.
  \item Find $\avg{x_n}$ and the variance of $x_n$.
  \item The drift $\avg{x_n}$ grows like $n$, the spread like $\sqrt n$.
        For $p=0.6$, after how many steps does the drift exceed the
        standard deviation?
\end{parts}

\solution

\solpart{a}{One step}
\begin{equation}
  \avg{s}=a\,p-a\,q=a(2p-1),\qquad \avg{s^2}=a^2 ,
\end{equation}
so
\begin{equation}
  \operatorname{Var}s=\avg{s^2}-\avg{s}^2=a^2\bigl[1-(2p-1)^2\bigr]
  =a^2\bigl[(1-2p+1)(1+2p-1)\bigr]=\ans{4a^2pq}.
\end{equation}
Check: $p=\tfrac12$ gives $a^2$, the unbiased walk; $p=1$ gives $0$, a
deterministic walk.

\solpart{b}{After $n$ steps}
Means add for any steps; variances add because the steps are independent:
\begin{equation}
  \ans{\avg{x_n}=na(2p-1),\qquad \operatorname{Var}x_n=4na^2pq}.
\end{equation}

\solpart{c}{Drift against spread}
The drift exceeds the standard deviation when
\begin{equation}
  na\,\lvert 2p-1\rvert>2a\sqrt{npq}
  \iff n>\frac{4pq}{(2p-1)^2}.
\end{equation}
For $p=0.6$: $4pq=0.96$ and $(2p-1)^2=0.04$, so $n>24$:
\ans{\text{from the 25th step on}}.

%% ----------------------------------------------------------------------------
\problem{How long does diffusion take?}{New}{\S2}

A small molecule in water has $D=1.0\times10^{-9}\,\mathrm{m^2\,s^{-1}}$.
Use the one-dimensional result $\avg{x^2}=2Dt$.
\begin{parts}
  \item How long until the root-mean-square displacement is $1\,\mathrm{mm}$?
  \item How long does it take to cross a cell, $10\,\mu\mathrm{m}$?
  \item Explain the ratio of the two times in one line.
\end{parts}

\solution

\solpart{a}{One millimetre}
\begin{equation}
  t=\frac{\avg{x^2}}{2D}=\frac{(10^{-3}\,\mathrm m)^2}{2\times10^{-9}\,\mathrm{m^2\,s^{-1}}}
  =\ans{500\,\mathrm s}\approx 8\,\mathrm{min}.
\end{equation}

\solpart{b}{Ten micrometres}
\begin{equation}
  t=\frac{(10^{-5}\,\mathrm m)^2}{2\times10^{-9}\,\mathrm{m^2\,s^{-1}}}
  =\ans{0.05\,\mathrm s}.
\end{equation}

\solpart{c}{The ratio}
The time grows as the square of the distance: $100$ times the distance takes
$100^2=10^4$ times as long, and $500\,\mathrm s/0.05\,\mathrm s=10^4$.

%% ============================================================================
\probchapter{Steady states}{Notes \S3}

%% ----------------------------------------------------------------------------
\problem{Flux through a membrane}{New}{\S3}

A membrane of thickness $L$ separates two well-stirred solutions. Inside the
membrane the concentration obeys the diffusion equation with
$c(0,t)=c_0$ and $c(L,t)=c_L$ held fixed.
\begin{parts}
  \item Find the steady profile $c(x)$.
  \item Find the steady flux $J=-D\,\partial c/\partial x$.
  \item Evaluate $J$ for $D=1.0\times10^{-9}\,\mathrm{m^2\,s^{-1}}$,
        $L=100\,\mu\mathrm m$, $c_0=1\,\mathrm{mol\,m^{-3}}$, $c_L=0$.
\end{parts}

\solution

\solpart{a}{The profile}
In a steady state $\partial c/\partial t=0$, so the diffusion equation leaves
$\d^2c/\d x^2=0$: the profile is a straight line. Fitting the two boundary
values,
\begin{equation}
  \ans{c(x)=c_0+(c_L-c_0)\,\frac xL}.
\end{equation}

\solpart{b}{The flux}
\begin{equation}
  J=-D\,\dv{c}{x}=\ans{D\,\frac{c_0-c_L}{L}},
\end{equation}
the same at every $x$, as it must be in a steady state (no accumulation
anywhere). It flows from high to low concentration.

\solpart{c}{Numbers}
\begin{equation}
  J=\frac{(1.0\times10^{-9}\,\mathrm{m^2\,s^{-1}})(1\,\mathrm{mol\,m^{-3}})}{1.0\times10^{-4}\,\mathrm m}
  =\ans{1.0\times10^{-5}\,\mathrm{mol\,m^{-2}\,s^{-1}}}.
\end{equation}


\end{document}
